\chapter{Natural Transformations}

\begin{overviewbox}
Chapter~3 built one layer of structure on top of categories: between any two
categories there are functors, and these functors themselves form a category.

This chapter adds a third layer. Between any two functors there are
\emph{natural transformations} --- morphisms between morphisms. A natural
transformation is a coordinated way of moving from one functor to another,
where ``coordinated'' has a very precise meaning: a single commutative square
must hold for every morphism in the source category.

By the end of this chapter we will have built a three-level hierarchy:
objects, morphisms, transformations between morphisms. Together with the
functor categories that hold them, these three levels are the structural
backbone of every result in the rest of the book.
\end{overviewbox}


% ═══════════════════════════════════════════════════════════════════════════
\section{What a Natural Transformation Is}
\label{sec:nat-trans}
% ═══════════════════════════════════════════════════════════════════════════

\subsection{Informally: A Coordinated Family of Morphisms}

Suppose we have two functors $F$ and $G$, both going from category
$\mathcal{C}$ to category $\mathcal{D}$. They send each object $A \in
\mathcal{C}$ to a (possibly different) object in $\mathcal{D}$. A natural
transformation from $F$ to $G$ is a recipe that, at every object $A$ of
$\mathcal{C}$, gives a morphism in $\mathcal{D}$ from $F(A)$ to $G(A)$.

But these morphisms cannot be chosen independently. They must be
\emph{coordinated} so that they agree with how $F$ and $G$ act on
morphisms. That coordination requirement is the entire content of the
definition.

\subsection{Expanding: Two Functors, a Family of Bridges}

\begin{example}{Two Recipes, One Translation}
Suppose the recipe category $\mathcal{R}$ (Chop $\to$ Sauté $\to$ Add) maps
to the approval chain $\mathcal{A}$ (Request $\to$ Approved $\to$ Final) in
two different ways:
\begin{itemize}
  \item Functor $F$: every step maps to Request.
  \item Functor $G$: Chop $\mapsto$ Request, Sauté $\mapsto$ Approved,
        Add $\mapsto$ Final.
\end{itemize}

A natural transformation $\alpha : F \Rightarrow G$ supplies, for each recipe
step $A$, a morphism in $\mathcal{A}$ from $F(A)$ to $G(A)$:
\begin{align*}
  \alpha_{\mathrm{Chop}} &: \mathrm{Request} \to \mathrm{Request} \\
  \alpha_{\mathrm{Saut\acute{e}}} &: \mathrm{Request} \to \mathrm{Approved} \\
  \alpha_{\mathrm{Add}}  &: \mathrm{Request} \to \mathrm{Final}
\end{align*}

In a thin category like $\mathcal{A}$, each of these is the unique morphism
between its endpoints. The natural transformation $\alpha$ provides
``advancement arrows'' that move from $F$'s constant Request view to $G$'s
finer-grained view, one step at a time.
\end{example}

\begin{dialog}
When I encounter $\alpha : F \Rightarrow G$, I picture two parallel functor
arrows from $\mathcal{C}$ to $\mathcal{D}$ and a sheaf of small morphisms
in $\mathcal{D}$ connecting their object-by-object outputs:
\medskip
\begin{center}
\begin{tikzcd}[column sep=4em, row sep=0em]
  \mathcal{C}
    \arrow[r, bend left=40, "F"{name=U}]
    \arrow[r, bend right=40, "G"{name=D, swap}]
  & \mathcal{D}
  \arrow[from=U, to=D, Rightarrow, "\alpha"]
\end{tikzcd}
\end{center}
\medskip
The double arrow $\Rightarrow$ is the visual signature of a natural
transformation. It lives ``between'' two functors.
\end{dialog}

\subsection{Formalizing: The Naturality Condition}

Let $F, G : \mathcal{C} \to \mathcal{D}$ be functors. A \term{natural
transformation} $\alpha : F \Rightarrow G$ consists of:

\begin{itemize}
  \item For each object $A \in \mathrm{Ob}(\mathcal{C})$, a morphism
        \(\alpha_A : F(A) \to G(A)\) in $\mathcal{D}$, called the
        \term{component of $\alpha$ at $A$}.
\end{itemize}

Subject to the \term{naturality condition}: for every morphism
$f : A \to B$ in $\mathcal{C}$, the following diagram commutes:
\medskip
\begin{center}
\begin{tikzcd}[column sep=4em, row sep=3em]
  F(A) \arrow[r, "\alpha_A"] \arrow[d, "F(f)"']
  & G(A) \arrow[d, "G(f)"] \\
  F(B) \arrow[r, "\alpha_B"']
  & G(B)
\end{tikzcd}
\end{center}
\medskip
That is:
\begin{equation*}
  G(f) \circ \alpha_A \;=\; \alpha_B \circ F(f).
\end{equation*}

This square is called the \term{naturality square} at $f$.

\begin{howtosay}
\begin{itemize}
  \item $\alpha : F \Rightarrow G$ is read ``alpha, a natural transformation
        from $F$ to $G$'' or ``$F$ goes to $G$ via $\alpha$.'' The double
        arrow is often called \emph{double-arrow}, \emph{nat-arrow}, or
        \emph{2-arrow}.
  \item $\alpha_A$ is ``the component of $\alpha$ at $A$,'' or
        ``alpha sub-$A$,'' or simply ``alpha at $A$.''
  \item Some authors write $\alpha : F \to G$ when the context is clearly
        functors. Either notation is standard.
  \item The naturality square ``commutes'' or, less formally, ``is natural
        in $A$'' (when we want to emphasize that this commutativity must hold
        as $A$ varies).
\end{itemize}
\end{howtosay}

\subsection{Reorganizing: Verifying Naturality on Our Example}

Return to $\alpha : F \Rightarrow G$ from the recipe example. For naturality
to hold, every morphism $f : A \to B$ in $\mathcal{R}$ must produce a
commutative square. Take $f : \mathrm{Chop} \to \mathrm{Saut\acute{e}}$.
The square is:
\begin{center}
\begin{tikzcd}[column sep=5em, row sep=3em]
  F(\mathrm{Chop}) = \mathrm{Request}
    \arrow[r, "\alpha_{\mathrm{Chop}}"]
    \arrow[d, "F(f) = \mathrm{id}_{\mathrm{Request}}"']
  & G(\mathrm{Chop}) = \mathrm{Request}
    \arrow[d, "G(f)"] \\
  F(\mathrm{Saut\acute{e}}) = \mathrm{Request}
    \arrow[r, "\alpha_{\mathrm{Saut\acute{e}}}"']
  & G(\mathrm{Saut\acute{e}}) = \mathrm{Approved}
\end{tikzcd}
\end{center}
Both paths from $\mathrm{Request}$ (top-left) to $\mathrm{Approved}$
(bottom-right) compose to the unique morphism between those two objects in
$\mathcal{A}$. The square commutes because $\mathcal{A}$ is thin --- there
is only one possible composite, so there is nothing to fail. The other
naturality squares (for $g$, $g \circ f$, and the identities) commute for
the same reason.

The natural transformation $\alpha$ is valid.

\subsection{Simplifying: Naturality as One Equation}

For every $f : A \to B$ in $\mathcal{C}$:
\begin{equation*}
  G(f) \circ \alpha_A \;=\; \alpha_B \circ F(f).
\end{equation*}

That equation, quantified over all $f$, is the entire definition of
naturality. Everything else --- the components, the diagrams, the names ---
is bookkeeping around this single equation.

\subsection{Formally: Notation and Conventions}

\formalnote{Notational summary}
For $F, G : \mathcal{C} \to \mathcal{D}$ and $\alpha : F \Rightarrow G$:
\begin{align*}
  \alpha_A &: F(A) \to G(A)
    && \text{(component at $A$)} \\
  G(f) \circ \alpha_A &= \alpha_B \circ F(f)
    && \text{(naturality at $f : A \to B$)}
\end{align*}

\formalnote{Identity natural transformation}
For any functor $F : \mathcal{C} \to \mathcal{D}$, there is an identity
natural transformation $\mathrm{id}_F : F \Rightarrow F$ whose components are
the identities of $\mathcal{D}$:
\begin{equation*}
  (\mathrm{id}_F)_A \;=\; \mathrm{id}_{F(A)}.
\end{equation*}
Naturality is automatic: $F(f) \circ \mathrm{id}_{F(A)} = F(f) =
\mathrm{id}_{F(B)} \circ F(f)$ by the unit laws.

\formalnote{Components determine the transformation}
A natural transformation is fully specified by its components together with
the verification of naturality. Two natural transformations $\alpha, \beta :
F \Rightarrow G$ are equal iff $\alpha_A = \beta_A$ for every $A$. There is
no extra data beyond the components.


% ═══════════════════════════════════════════════════════════════════════════
\section{Vertical Composition}
\label{sec:vertical-comp}
% ═══════════════════════════════════════════════════════════════════════════

\subsection{Informally: Stacking Transformations}

If $\alpha : F \Rightarrow G$ takes us from $F$ to $G$, and $\beta : G
\Rightarrow H$ takes us from $G$ to $H$, then doing $\alpha$ followed by
$\beta$ should take us from $F$ to $H$. This composed transformation is
denoted $\beta \cdot \alpha$ (or sometimes just $\beta \alpha$) and is called
the \term{vertical composition}.

It is called \emph{vertical} because in the standard diagram, we picture the
two transformations stacked top-to-bottom, with the composed transformation
running the full height.

\subsection{Expanding: Stacking Components}

\begin{center}
\begin{tikzcd}[column sep=5em, row sep=0em]
  \mathcal{C}
    \arrow[r, bend left=60, "F"{name=F1}]
    \arrow[r, "G"{name=G1, description}]
    \arrow[r, bend right=60, "H"{name=H1, swap}]
  & \mathcal{D}
  \arrow[from=F1, to=G1, Rightarrow, "\alpha"]
  \arrow[from=G1, to=H1, Rightarrow, "\beta"]
\end{tikzcd}
\end{center}

At each object $A$ of $\mathcal{C}$, we have
\begin{align*}
  \alpha_A &: F(A) \to G(A), \\
  \beta_A  &: G(A) \to H(A).
\end{align*}
These two morphisms compose in $\mathcal{D}$:
\begin{equation*}
  (\beta \cdot \alpha)_A \;\coloneqq\; \beta_A \circ \alpha_A
  \;:\; F(A) \to H(A).
\end{equation*}

The component of the composed transformation at $A$ is the composite of
the two original components at $A$.

\subsection{Formalizing: Naturality of the Composite}

We must check that $\beta \cdot \alpha$ is itself a natural transformation
--- i.e., that its components $(\beta \cdot \alpha)_A$ satisfy the
naturality condition. For any $f : A \to B$ in $\mathcal{C}$:
\begin{align*}
  H(f) \circ (\beta \cdot \alpha)_A
    &= H(f) \circ \beta_A \circ \alpha_A
       && \text{(definition of composite)} \\
    &= \beta_B \circ G(f) \circ \alpha_A
       && \text{(naturality of $\beta$ at $f$)} \\
    &= \beta_B \circ \alpha_B \circ F(f)
       && \text{(naturality of $\alpha$ at $f$)} \\
    &= (\beta \cdot \alpha)_B \circ F(f)
       && \text{(definition of composite)}.
\end{align*}
The two outer expressions equal, so the naturality square for
$\beta \cdot \alpha$ at $f$ commutes. $\beta \cdot \alpha$ is a natural
transformation.

\begin{howtosay}
\begin{itemize}
  \item $\beta \cdot \alpha$ is ``beta vertical alpha'' or ``beta after
        alpha'' or simply ``beta composed with alpha.''
  \item The dot $\cdot$ is the conventional symbol for vertical composition.
        Some authors drop it entirely and write $\beta \alpha$.
  \item Order: $\beta \cdot \alpha$ does $\alpha$ first, then $\beta$ ---
        right to left, matching morphism composition.
\end{itemize}
\end{howtosay}

\subsection{Reorganizing: Associativity and Units}

Vertical composition is associative: for $\alpha : F \Rightarrow G$,
$\beta : G \Rightarrow H$, $\gamma : H \Rightarrow K$:
\begin{equation*}
  (\gamma \cdot \beta) \cdot \alpha \;=\; \gamma \cdot (\beta \cdot \alpha).
\end{equation*}
This follows componentwise from the associativity of composition in
$\mathcal{D}$:
\begin{align*}
  ((\gamma \cdot \beta) \cdot \alpha)_A
    &= (\gamma \cdot \beta)_A \circ \alpha_A \\
    &= (\gamma_A \circ \beta_A) \circ \alpha_A \\
    &= \gamma_A \circ (\beta_A \circ \alpha_A)
       && \text{(associativity in $\mathcal{D}$)} \\
    &= \gamma_A \circ (\beta \cdot \alpha)_A \\
    &= (\gamma \cdot (\beta \cdot \alpha))_A.
\end{align*}

Identity transformations are units for vertical composition:
\begin{equation*}
  \alpha \cdot \mathrm{id}_F \;=\; \alpha
  \;=\; \mathrm{id}_G \cdot \alpha
  \qquad \text{for } \alpha : F \Rightarrow G.
\end{equation*}

\subsection{Simplifying: A Familiar Pattern}

We have:
\begin{itemize}
  \item Objects: functors $F, G, H, \ldots$ all $\mathcal{C} \to \mathcal{D}$.
  \item Morphisms: natural transformations $\alpha : F \Rightarrow G$.
  \item Composition: vertical composition $\beta \cdot \alpha$, associative.
  \item Identity: $\mathrm{id}_F : F \Rightarrow F$, neutral.
\end{itemize}

The pattern is exactly that of a category. We name the resulting category in
the next section.

\subsection{Formally: Componentwise Verification Pattern}

\formalnote{The componentwise pattern}
Most facts about natural transformations are proved by ``checking
componentwise'': reduce an equation between natural transformations to an
equation between their components at each object, then verify the latter
using the rules of the target category $\mathcal{D}$. Naturality of the
candidate composite or identity is then verified separately.

The proof above of associativity is a paradigmatic example. The reader
should expect to recognize this pattern often.


% ═══════════════════════════════════════════════════════════════════════════
\section{Functor Categories}
\label{sec:functor-cats}
% ═══════════════════════════════════════════════════════════════════════════

\subsection{Informally: Functors as Objects}

Take any two categories $\mathcal{C}$ and $\mathcal{D}$. The functors from
$\mathcal{C}$ to $\mathcal{D}$ are objects in a new category, whose
morphisms are the natural transformations between them. This new category
is called the \term{functor category} from $\mathcal{C}$ to $\mathcal{D}$.

\subsection{Expanding: Putting the Pieces Together}

We have already collected the ingredients in the previous section. We name
them once more, now framed as a category:

\begin{itemize}
  \item \textbf{Objects}: functors $F : \mathcal{C} \to \mathcal{D}$.
  \item \textbf{Morphisms}: natural transformations $\alpha : F \Rightarrow G$
        between such functors.
  \item \textbf{Composition}: vertical composition $\beta \cdot \alpha$.
  \item \textbf{Identity}: identity natural transformation $\mathrm{id}_F$.
\end{itemize}

The axioms have been verified:
\begin{align*}
  (\gamma \cdot \beta) \cdot \alpha
    &= \gamma \cdot (\beta \cdot \alpha)
    && \text{(associativity)} \\
  \alpha \cdot \mathrm{id}_F
    &= \alpha
    = \mathrm{id}_G \cdot \alpha
    && \text{(unit laws)}.
\end{align*}

This is a category.

\subsection{Formalizing: Notation}

The functor category from $\mathcal{C}$ to $\mathcal{D}$ is denoted in
several equivalent ways:
\begin{equation*}
  [\mathcal{C}, \mathcal{D}]
  \;=\; \mathcal{D}^{\mathcal{C}}
  \;=\; \mathrm{Fun}(\mathcal{C}, \mathcal{D}).
\end{equation*}

\begin{howtosay}
\begin{itemize}
  \item $[\mathcal{C}, \mathcal{D}]$ — ``the functor category from
        $\mathcal{C}$ to $\mathcal{D}$,'' or ``brackets $\mathcal{C}$,
        $\mathcal{D}$.'' Common in modern texts.
  \item $\mathcal{D}^{\mathcal{C}}$ — ``$\mathcal{D}$ to the $\mathcal{C}$''
        or ``$\mathcal{D}$ exponential $\mathcal{C}$.'' This notation
        suggests an analogy with exponentiation that becomes precise once
        we have products.
  \item $\mathrm{Fun}(\mathcal{C}, \mathcal{D})$ — ``Fun of $\mathcal{C}$,
        $\mathcal{D}$'' or simply ``the functors from $\mathcal{C}$ to
        $\mathcal{D}$.''
\end{itemize}
This book uses $[\mathcal{C}, \mathcal{D}]$ by default and falls back to
$\mathcal{D}^{\mathcal{C}}$ when the exponential analogy is the point.
\end{howtosay}

\subsection{Reorganizing: A Familiar Example}

The category $[\mathbf{1}, \mathcal{C}]$ has:
\begin{itemize}
  \item Objects: functors from $\mathbf{1}$ to $\mathcal{C}$, which (we
        showed in \S3.2) correspond exactly to objects of $\mathcal{C}$.
  \item Morphisms: natural transformations between such functors. A natural
        transformation between two such functors $F$ and $G$, picking out
        objects $A = F(*)$ and $B = G(*)$, has one component
        $\alpha_* : A \to B$. Naturality at $\mathrm{id}_*$ is automatic.
\end{itemize}

So $[\mathbf{1}, \mathcal{C}]$ is essentially the same as $\mathcal{C}$
itself --- objects are objects of $\mathcal{C}$, morphisms are morphisms of
$\mathcal{C}$. This is our first ``equivalence of categories,'' though we
have not yet defined that term precisely. It will become precise in a later
chapter.

\subsection{Simplifying: The Exponential Notation}

\begin{equation*}
  \mathcal{D}^{\mathcal{C}}
  \;=\; \text{the category whose objects are functors}\; \mathcal{C} \to \mathcal{D}.
\end{equation*}

\subsection{Formally: Functor Categories Are Categories}

\formalnote{Existence}
For any (small) $\mathcal{C}$ and any $\mathcal{D}$, the functor category
$[\mathcal{C}, \mathcal{D}]$ exists and is itself a category.

\formalnote{Size considerations}
If $\mathcal{C}$ is small and $\mathcal{D}$ is locally small, then
$[\mathcal{C}, \mathcal{D}]$ is locally small. If both are small,
$[\mathcal{C}, \mathcal{D}]$ is small.

\formalnote{Why this matters}
Functor categories are the natural setting for many results that follow.
The Yoneda lemma, presheaves, sheaves, and most universal constructions
either live inside a functor category or are characterized in terms of one.


% ═══════════════════════════════════════════════════════════════════════════
\section{Natural Isomorphism}
\label{sec:nat-iso}
% ═══════════════════════════════════════════════════════════════════════════

\subsection{Informally: When Two Functors Are ``the Same''}

In Chapter~2 we said two objects $A$ and $B$ of a category are
\emph{isomorphic} when there is an invertible morphism between them. The
same idea, applied to the functor category $[\mathcal{C}, \mathcal{D}]$,
gives a notion of two \emph{functors} being isomorphic. Two functors $F$ and
$G$ are \emph{naturally isomorphic} when there is an invertible natural
transformation between them.

A natural isomorphism is just an isomorphism in a functor category, no more,
no less. But because functors carry so much structure, this is one of the
most useful notions in the entire subject.

\subsection{Expanding: Componentwise Invertibility}

A natural transformation $\alpha : F \Rightarrow G$ has an inverse
$\alpha^{-1} : G \Rightarrow F$ exactly when there is a natural
transformation in the other direction such that the two vertical composites
are identity natural transformations:
\begin{align*}
  \alpha^{-1} \cdot \alpha &= \mathrm{id}_F, \\
  \alpha \cdot \alpha^{-1} &= \mathrm{id}_G.
\end{align*}

By the componentwise pattern, these equations hold iff for every $A$:
\begin{align*}
  (\alpha^{-1})_A \circ \alpha_A &= \mathrm{id}_{F(A)}, \\
  \alpha_A \circ (\alpha^{-1})_A &= \mathrm{id}_{G(A)}.
\end{align*}

So $\alpha$ is a natural isomorphism iff every component $\alpha_A$ is an
isomorphism in $\mathcal{D}$, with $(\alpha^{-1})_A = (\alpha_A)^{-1}$.

\subsection{Formalizing: Definition and Notation}

A natural transformation $\alpha : F \Rightarrow G$ is a \term{natural
isomorphism} if there exists a natural transformation $\alpha^{-1} :
G \Rightarrow F$ with $\alpha^{-1} \cdot \alpha = \mathrm{id}_F$ and
$\alpha \cdot \alpha^{-1} = \mathrm{id}_G$.

Equivalently: $\alpha$ is a natural isomorphism iff every component
$\alpha_A : F(A) \to G(A)$ is an isomorphism in $\mathcal{D}$.

Two functors $F$ and $G$ are \term{naturally isomorphic}, written
$F \cong G$, if there is a natural isomorphism between them.

\begin{howtosay}
\begin{itemize}
  \item $F \cong G$ — ``$F$ is naturally isomorphic to $G$,'' or ``$F$ and
        $G$ are naturally isomorphic,'' or (in informal speech) ``$F$ and
        $G$ are the same up to natural isomorphism.''
  \item The double arrow with a tilde, $F \xRightarrow{\sim} G$ or
        $F \overset{\sim}{\Rightarrow} G$, indicates that the natural
        transformation $\alpha$ is a natural isomorphism. Reads as
        ``alpha is a natural iso from $F$ to $G$.''
  \item Some authors say ``isomorphic'' (without ``naturally'') when the
        context is clear that the isomorphism is in a functor category.
\end{itemize}
\end{howtosay}

\subsection{Reorganizing: Naturality of the Inverse Is Automatic}

It is worth checking explicitly that if $\alpha$ has componentwise inverses
$(\alpha_A)^{-1}$ in $\mathcal{D}$, then the family of inverses is itself a
natural transformation. We verify the naturality square for the candidate
inverse:
\begin{align*}
  F(f) \circ (\alpha_A)^{-1}
    &= (\alpha_B)^{-1} \circ \alpha_B \circ F(f) \circ (\alpha_A)^{-1}
       && \text{(inserting identity)} \\
    &= (\alpha_B)^{-1} \circ G(f) \circ \alpha_A \circ (\alpha_A)^{-1}
       && \text{(naturality of $\alpha$)} \\
    &= (\alpha_B)^{-1} \circ G(f)
       && \text{(inverse cancels)}.
\end{align*}
The outer equality is precisely the naturality square for the inverse at
$f : A \to B$. So having componentwise inverses suffices --- the naturality
of the inverse comes for free.

\subsection{Simplifying: Three Equivalent Statements}

The following are equivalent for $\alpha : F \Rightarrow G$:
\begin{align*}
  &\text{(i)} \quad \alpha \text{ is an isomorphism in } [\mathcal{C}, \mathcal{D}]. \\
  &\text{(ii)} \quad \text{there exists } \alpha^{-1} \text{ with }
    \alpha^{-1} \cdot \alpha = \mathrm{id}_F \text{ and }
    \alpha \cdot \alpha^{-1} = \mathrm{id}_G. \\
  &\text{(iii)} \quad \text{every component } \alpha_A \text{ is an isomorphism in } \mathcal{D}.
\end{align*}

The third is by far the most useful test: to check a natural isomorphism,
check that every component is an isomorphism in the target category.

\subsection{Formally: Natural Isomorphism in Practice}

\formalnote{Natural vs.\ unnatural isomorphism}
Sometimes one can find isomorphisms $F(A) \cong G(A)$ for every object
$A$, but no single choice of isomorphisms that is natural in $A$. In that
case $F$ and $G$ are \emph{not} naturally isomorphic, even though they are
``pointwise isomorphic.'' Naturality forces the family of isomorphisms to
cohere with morphisms.

This distinction is one of the historical motivations for the entire
subject: Eilenberg and Mac~Lane introduced natural transformations
precisely to formalize ``natural'' versus ``arbitrary'' choices.

\formalnote{Functor categories have isomorphism}
Just as in any category, ``naturally isomorphic'' is an equivalence
relation on the objects of $[\mathcal{C}, \mathcal{D}]$ (the functors). The
proof is the same proof as for ordinary isomorphism --- it works in any
category, and the functor category is no exception.


% ═══════════════════════════════════════════════════════════════════════════
\section*{Exercises}
\addcontentsline{toc}{section}{Exercises}
% ═══════════════════════════════════════════════════════════════════════════

\begin{exercisesbox}
\begin{qa}
\qa{Functors are maps between categories. What is a map between functors?}{A natural transformation.}
\qa{If $F, G : \mathcal{C} \to \mathcal{D}$, what does $\alpha : F \Rightarrow G$ consist of?}{A morphism $\alpha_A : F(A) \to G(A)$ in $\mathcal{D}$ for every object $A$ of $\mathcal{C}$.}
\qa{What is $\alpha_A$ called?}{The component of $\alpha$ at $A$.}
\qa{What constraint do the components satisfy?}{The naturality square: for every $f : A \to B$, $G(f) \circ \alpha_A = \alpha_B \circ F(f)$.}
\qa{What does the naturality square look like as a commuting diagram?}{A square with $F(A), G(A), F(B), G(B)$ at the corners, and $F(f), G(f), \alpha_A, \alpha_B$ on the sides.}
\qa{Why is it called \emph{natural}?}{Because the components are chosen in a way that respects all morphisms, not just one.}
\qa{What is the identity natural transformation on $F$?}{$\mathrm{id}_F$, with $(\mathrm{id}_F)_A = \mathrm{id}_{F(A)}$.}
\qa{Does $\mathrm{id}_F$ satisfy naturality?}{Yes, trivially: both paths around the square reduce to $F(f)$.}
\qa{What is vertical composition?}{Composing $\alpha : F \Rightarrow G$ with $\beta : G \Rightarrow H$ to get $\beta \cdot \alpha : F \Rightarrow H$.}
\qa{What is the component formula for $\beta \cdot \alpha$?}{$(\beta \cdot \alpha)_A = \beta_A \circ \alpha_A$.}
\qa{Is vertical composition associative?}{Yes.}
\qa{What is the functor category $[\mathcal{C}, \mathcal{D}]$?}{The category whose objects are functors $\mathcal{C} \to \mathcal{D}$ and whose morphisms are natural transformations between them.}
\qa{What is the identity morphism in $[\mathcal{C}, \mathcal{D}]$?}{The identity natural transformation $\mathrm{id}_F$.}
\qa{What are two other notations for $[\mathcal{C}, \mathcal{D}]$?}{$\mathcal{D}^{\mathcal{C}}$ and $\mathrm{Fun}(\mathcal{C}, \mathcal{D})$.}
\qa{When is $\alpha : F \Rightarrow G$ called a natural isomorphism?}{When every component $\alpha_A$ is an isomorphism in $\mathcal{D}$.}
\qa{If every component is an iso, is the inverse natural too?}{Yes. The pointwise inverses assemble into a natural transformation.}
\qa{What does $F \cong G$ mean in $[\mathcal{C}, \mathcal{D}]$?}{There exists a natural isomorphism $F \Rightarrow G$.}
\qa{Can $F$ and $G$ act differently on objects yet be naturally isomorphic?}{Yes; their object-values need only be isomorphic, not equal.}
\qa{Why is naturality the right notion of ``map of functors''?}{Without it, components could behave incoherently across morphisms; naturality enforces coherence.}
\qa{What is horizontal composition?}{Composing $\alpha : F \Rightarrow G$ in $[\mathcal{C}, \mathcal{D}]$ with $\beta : H \Rightarrow K$ in $[\mathcal{D}, \mathcal{E}]$ to get $\beta \ast \alpha : H \circ F \Rightarrow K \circ G$.}
\qa{What is the whiskering $\beta F$?}{Horizontal composition $\beta \ast \mathrm{id}_F$: precomposing each component of $\beta$ with $F$.}
\qa{What is the component of $\beta F$ at $A$?}{$\beta_{F(A)}$.}
\qa{What is $\alpha_A : F(A) \to G(A)$ best read as?}{``$\alpha$ component at $A$,'' a single morphism in $\mathcal{D}$.}
\qa{Is naturality a property or extra data?}{A property: the components are the data, and naturality is a condition they must satisfy.}
\qa{Do natural transformations between functors $\mathcal{C} \to \mathcal{D}$ form a set?}{They form a hom-set in $[\mathcal{C}, \mathcal{D}]$.}
\qa{Why do we need naturality if pointwise isomorphism would do?}{Pointwise iso ignores coherence; many ``isomorphic'' functors fail to be naturally isomorphic.}
\qa{Where does the symbol $\Rightarrow$ come from?}{It distinguishes a 2-cell (between 1-cells) from a 1-cell (between 0-cells).}
\qa{Is $F \Rightarrow G$ a different thing from $F \to G$?}{Conventionally yes: $\to$ is a morphism, $\Rightarrow$ is a natural transformation between functors.}
\qa{What is the simplest non-trivial natural transformation example?}{The identity-to-identity isomorphism, or the inclusion of a constant functor into a non-constant one.}
\qa{Why did we want functor categories at all?}{So that we can apply categorical reasoning \emph{to functors themselves} --- the same machinery, one level up.}
\end{qa}
\end{exercisesbox}


% ═══════════════════════════════════════════════════════════════════════════
\section*{Chapter Summary}
\addcontentsline{toc}{section}{Chapter Summary}
% ═══════════════════════════════════════════════════════════════════════════

\begin{summarybox}
A \textbf{natural transformation} $\alpha : F \Rightarrow G$ between two
parallel functors $F, G : \mathcal{C} \to \mathcal{D}$ consists of a
component $\alpha_A : F(A) \to G(A)$ for each object $A$, such that for
every morphism $f : A \to B$ in $\mathcal{C}$, the naturality square
commutes:
\begin{equation*}
  G(f) \circ \alpha_A \;=\; \alpha_B \circ F(f).
\end{equation*}

\textbf{Vertical composition}: $(\beta \cdot \alpha)_A = \beta_A \circ
\alpha_A$. Associative; identity natural transformations are units.

\textbf{Functor categories}: $[\mathcal{C}, \mathcal{D}]$ (also written
$\mathcal{D}^{\mathcal{C}}$) has functors as objects and natural
transformations as morphisms.

\textbf{Natural isomorphism} $F \cong G$: an isomorphism in
$[\mathcal{C}, \mathcal{D}]$, equivalently a natural transformation whose
every component is an isomorphism in $\mathcal{D}$.

\medskip
\textbf{The three levels:}
\begin{enumerate}
  \item objects (live in a category),
  \item morphisms (between objects in the same category),
  \item natural transformations (between functors with the same source and
        target categories).
\end{enumerate}
\end{summarybox}

% ═══════════════════════════════════════════════════════════════════════════
\section*{Formal Restatement}
\addcontentsline{toc}{section}{Formal Restatement}
% ═══════════════════════════════════════════════════════════════════════════

\begin{formalrestatement}
\textbf{Natural transformation.}\quad
For $F, G : \mathcal{C} \to \mathcal{D}$, a natural transformation
$\alpha : F \Rightarrow G$ assigns to each $A \in \mathcal{C}$ a morphism
$\alpha_A : F(A) \to G(A)$ in $\mathcal{D}$, subject to:
\begin{equation*}
  G(f) \circ \alpha_A \;=\; \alpha_B \circ F(f)
  \qquad \text{for all } f : A \to B \text{ in } \mathcal{C}.
\end{equation*}

\medskip
\textbf{Identity.}\quad $\mathrm{id}_F : F \Rightarrow F$ has components
$(\mathrm{id}_F)_A = \mathrm{id}_{F(A)}$.

\medskip
\textbf{Vertical composition.}\quad For $\alpha : F \Rightarrow G$ and
$\beta : G \Rightarrow H$:
\begin{equation*}
  (\beta \cdot \alpha)_A \;=\; \beta_A \circ \alpha_A.
\end{equation*}
Associative; $\mathrm{id}_G \cdot \alpha = \alpha = \alpha \cdot \mathrm{id}_F$.

\medskip
\textbf{Functor category.}\quad
$[\mathcal{C}, \mathcal{D}] = \mathcal{D}^{\mathcal{C}}$ has functors
$\mathcal{C} \to \mathcal{D}$ as objects and natural transformations as
morphisms.

\medskip
\textbf{Natural isomorphism.}\quad
$\alpha : F \Rightarrow G$ is a natural isomorphism iff $\alpha$ is
invertible in $[\mathcal{C}, \mathcal{D}]$ iff every component $\alpha_A$
is an isomorphism in $\mathcal{D}$. Two functors $F$ and $G$ are naturally
isomorphic, $F \cong G$, when such an $\alpha$ exists.

\medskip
\noindent\textit{Chapter~5 introduces universal properties, beginning with
initial and terminal objects: the first examples of objects characterized
by their relationship to everything else.}
\end{formalrestatement}
